\documentclass[11pt,a4paper]{article}

\usepackage[utf8]{inputenc}
\usepackage[T1]{fontenc}
\usepackage[spanish]{babel}
\usepackage{geometry}
\usepackage{booktabs}
\usepackage{longtable}
\usepackage{amsmath}
\usepackage{amssymb}
\usepackage{hyperref}
\usepackage{xcolor}
\usepackage{url}

\geometry{margin=2.4cm}
\hypersetup{colorlinks=true, linkcolor=blue, urlcolor=blue}
\renewcommand{\arraystretch}{1.12}

\title{\textbf{DeMemteAttractor E5: memoria atractora con compuerta anti-pareidolia para clasificación robusta}\\
\large Recapitulación metodológica de \texttt{e5\_final\_clean.ipynb}}
\author{Christian Borasino}
\date{}

\begin{document}
\maketitle

\begin{abstract}
Este informe recapitula, en formato de artículo técnico, el contenido del
notebook \texttt{experiments/atracctor/e5\_final\_clean.ipynb}. La libreta
consolida la variante final \textit{E5 combined dropout + OOD tau 1.50} de
DeMemteAttractor, un modelo que combina una ResNet18 congelada, una memoria
latente vector-cuantizada, un módulo atractor residual y una compuerta de
ambigüedad guiada por incertidumbre, familiaridad, conflicto y riesgo de fuera
de distribución. La contribución práctica de E5 no consiste en maximizar una
sola métrica, sino en satisfacer simultáneamente precisión limpia, robustez ante
corrupciones, ordenamiento funcional de la compuerta y control de intervenciones
dañinas. En los artefactos cargados por el notebook, E5 obtiene 0.7653 de
precisión limpia, la mejor precisión corrupta promedio reportada, 0.4313, y es
la única variante con \texttt{strict\_acceptance\_success=True}. Este resultado
se logra al evitar la saturación del gate crudo, cerrar la memoria ante ruido
gaussiano severo y mantener intervención moderada frente a corrupciones más
recuperables como blur y cutout.
\end{abstract}

\section{Introducción}

\texttt{e5\_final\_clean.ipynb} es una versión depurada para preservar y explicar
la variante final ganadora del proyecto DeMemteAttractor. Su propósito declarado
es conservar la variante \texttt{E5 combined dropout + OOD tau 1.50} y reunir lo
necesario para cargar el checkpoint final, revisar métricas guardadas, graficar
el comportamiento de la compuerta, recomputar la evaluación completa si se
requiere y reentrenar E5 de forma opcional.

El problema experimental es la clasificación de \texttt{Flowers102} bajo
condiciones limpias y corruptas. La hipótesis operacional codificada en el
notebook es que una memoria latente puede recuperar representaciones dañadas,
pero que dicha memoria no debe intervenir indiscriminadamente. Si una
perturbación es recuperable, la memoria puede actuar como mecanismo de
completamiento; si la entrada se aleja del dominio familiar, una intervención
fuerte puede producir pareidolia, es decir, una corrección que induce una clase
errónea. Por ello, el éxito se define mediante una evaluación compuesta y no por
precisión aislada.

La variante final se denomina internamente
\texttt{e5\_combined\_dropout\_ood\_tau\_150} y se presenta en los artefactos
como \textit{E5 combined dropout + OOD tau 1.50}. Esta recapitulación toma
exclusivamente la arquitectura, la configuración y los artefactos que el propio
notebook carga.

\section{Método}

DeMemteAttractor parte de una ResNet18 preentrenada en ImageNet y truncada antes
de sus capas finales. El backbone produce un mapa espacial de 512 canales que se
mantiene congelado. Sobre esta representación se construye una ruta latente con
tres elementos: un proyector $1\times1$ hacia 128 canales, un cuantizador
vectorial 2D con 1024 embeddings y un desproyector $1\times1$ de vuelta a 512
canales. En paralelo, un atractor residual aprende correcciones token a token en
el espacio latente.

La forma general del modelo puede resumirse como:
\[
x \xrightarrow{\text{ResNet18}} f
\xrightarrow{\text{projector}} z
\xrightarrow{\text{attractor}} z_{\text{completed}}.
\]
El atractor no reemplaza directamente a $z$; produce una corrección residual:
\[
\Delta = z_{\text{completed}} - \operatorname{stopgrad}(z).
\]
La representación final queda modulada por una compuerta escalar por muestra:
\[
z_{\text{final}} = z + g\Delta.
\]
Después, $z_{\text{final}}$ se desproyecta al espacio de 512 canales y se
clasifica con una cabeza lineal. Una cabeza auxiliar clasifica las features
originales del backbone y cumple dos funciones: proveer una referencia base y
alimentar la señal de incertidumbre usada por el gate.

\subsection{Memoria vector-cuantizada}

El módulo \texttt{VectorQuantizer2D} transforma el mapa latente en tokens
espaciales y asigna cada token al embedding más cercano de un codebook de 1024
vectores. La distancia usada es cuadrática:
\[
d(z,e)=\|z\|^2 - 2ze^\top + \|e\|^2.
\]
El módulo devuelve el cuantizado con estimador straight-through, una pérdida VQ
y un mapa de discrepancia \texttt{dq\_map}. La pérdida VQ combina compromiso del
encoder y actualización del codebook:
\[
\mathcal{L}_{vq}
= c\,\operatorname{MSE}(z_e,\operatorname{stopgrad}(q))
+ \operatorname{MSE}(q,\operatorname{stopgrad}(z_e)),
\]
con $c=0.25$. El \texttt{dq\_map} se vuelve una señal central para juzgar
familiaridad: si la representación continua queda lejos del codebook, el modelo
interpreta que la entrada puede ser menos confiable para intervención
memorística.

\subsection{Atractor residual}

\texttt{AttractorMemory} es un MLP residual aplicado a cada token latente. Su
arquitectura es $128 \rightarrow 512 \rightarrow 128$ con activación
\texttt{GELU}. Al sumar la salida al token original, el atractor aprende
correcciones y no una sustitución completa de la representación. Esta decisión
es coherente con la función de E5: completar cuando conviene, pero conservar una
ruta directa disponible.

\subsection{Compuerta de ambigüedad}

\texttt{AmbiguityGate} recibe cuatro señales por muestra:
\[
\text{uncertainty},\quad \text{familiarity},\quad \text{conflict},\quad
1-\text{ood\_risk}.
\]
La incertidumbre es la entropía normalizada de los logits auxiliares. La
familiaridad es una gaussiana sobre el error VQ normalizado por estadísticas EMA:
\[
dq_{\text{norm}}=\frac{dq_{\text{score}}-\mu_{EMA}}{\sigma_{EMA}+10^{-5}},
\]
\[
\text{familiarity}
=\exp\left(
-\frac{(dq_{\text{norm}}-\text{midpoint})^2}{2\,\text{width}^2}
\right).
\]
El conflicto se obtiene como entropía normalizada de las asignaciones blandas del
cuantizador. El riesgo OOD se calcula con una sigmoide:
\[
\text{ood\_risk}=\sigma\left(\beta(dq_{\text{norm}}-\tau)\right).
\]

La compuerta final combina un prior interpretable con un MLP pequeño:
\[
\text{gate\_prior}
=\text{uncertainty}\cdot\text{familiarity}\cdot(1-\text{ood\_risk}),
\]
\[
g = floor + (1-floor)\cdot \text{gate\_prior}\cdot \text{gate\_raw}.
\]
El piso se fija en 0.02. Este diseño impide que el modelo dependa solo de una
red aprendida opaca: incluso si el MLP produce un \texttt{gate\_raw} alto, el
prior puede cerrar la intervención cuando la familiaridad cae o el riesgo OOD
sube.

\section{Entrenamiento}

El notebook conserva un entrenamiento en tres fases, aunque por defecto no lo
ejecuta porque \texttt{RUN\_TRAINING=False}. La presencia de este bloque permite
reproducir la construcción de E5 desde cero si se activa manualmente.

\subsection{Fase 1: estabilización de la memoria latente}

La primera fase congela backbone, atractor, gate y clasificadores. Solo entrena
proyector, cuantizador y desproyector. El objetivo es que la memoria latente
aprenda a reconstruir las features del backbone antes de recibir presión de
clasificación. Durante entrenamiento se enmascara una fracción de features con
\texttt{masked\_feature\_ratio=0.35}. La pérdida de fase 1 es:
\[
\mathcal{L}_{P1}
=0.5\mathcal{L}_{recon}
+0.25\mathcal{L}_{vq}
+\lambda_{sigreg}\mathcal{L}_{sigreg}.
\]
La regularización \texttt{SigReg} penaliza desviaciones de la covarianza latente
respecto a la identidad, con sketch de dimensión 64.

\subsection{Fase 2: aprendizaje de atractor, gate y cabezas}

La segunda fase congela la memoria latente ya inicializada y entrena el atractor,
la compuerta y las dos cabezas de clasificación. Cada lote se evalúa en versión
limpia y en versión corrupta. El objetivo latente limpio se calcula sin gradiente
y se usa como referencia de denoising para la entrada corrupta. La pérdida
incluye reconstrucción latente, clasificación limpia, entropía de gate y SigReg:
\[
\mathcal{L}_{P2}
=\mathcal{L}_{mse,dirty}
+0.3\mathcal{L}_{mse,clean}
+\mathcal{L}_{CE,clean}
+\lambda_H\mathcal{R}_{gate}
+\lambda_{raw}\mathcal{R}_{gate\_raw}
+\lambda_{sigreg}\mathcal{L}_{sigreg}.
\]
La función \texttt{phase2\_gate\_sanity} comprueba que el gate medio en un batch
corrupto no colapse fuera de $[0.1,0.9]$.

\subsection{Fase 3: ajuste conjunto y control anti-pareidolia}

La tercera fase optimiza clasificación limpia y corrupta, denoising, VQ,
regularización del gate, regularización latente y una pérdida anti-pareidolia.
En la configuración final de E5, \texttt{phase3\_memory\_grad\_mode=full}, por
lo que la memoria latente también puede ajustarse. La pérdida agregada es:
\[
\begin{aligned}
\mathcal{L}_{P3}=&\ 0.5(\mathcal{L}_{CE,clean}+\mathcal{L}_{CE,dirty})\\
&+0.5w_{denoise}(\mathcal{L}_{mse,clean}+\mathcal{L}_{mse,dirty})\\
&+0.5w_{vq}(\mathcal{L}_{vq,clean}+\mathcal{L}_{vq,dirty})\\
&+w_{anti}\mathcal{L}_{anti}
+\lambda_H\mathcal{R}_{gate}
+\lambda_{raw}\mathcal{R}_{gate\_raw}
+\lambda_{sigreg}\mathcal{L}_{sigreg}.
\end{aligned}
\]
La pérdida anti-pareidolia penaliza dos patrones: gate alto cuando la
familiaridad es baja, y gate alto cuando el clasificador auxiliar acierta pero la
salida final se equivoca. Este término alinea el entrenamiento con el criterio
de seguridad: la memoria no debe inventar estructura donde la señal de
familiaridad no la respalda.

\section{Definición de la variante E5}

El notebook define una única especificación final en \texttt{experiment\_specs}.
Todas las señales del gate permanecen activas:
\texttt{use\_uncertainty=True}, \texttt{use\_familiarity=True},
\texttt{use\_conflict=True} y \texttt{use\_ood=True}. Sobre la configuración base
se aplican los overrides de la Tabla~\ref{tab:e5_config}.

\begin{table}[h]
\centering
\begin{tabular}{ll}
\toprule
Parámetro & Valor en E5 \\
\midrule
\texttt{ood\_tau} & 1.5 \\
\texttt{ood\_beta} & 8.0 \\
\texttt{familiarity\_width} & 0.5 \\
\texttt{phase3\_lock\_familiarity} & \texttt{True} \\
\texttt{gate\_dropout} & 0.1 \\
\texttt{lr\_gate} & $1\cdot10^{-4}$ \\
\texttt{gate\_raw\_entropy\_reg} & 0.01 \\
\bottomrule
\end{tabular}
\caption{Modificaciones que definen la variante E5 respecto a la configuración base.}
\label{tab:e5_config}
\end{table}

Estos cambios atacan directamente el fallo observado en variantes anteriores: la
saturación del gate. Un \texttt{ood\_tau} de 1.5 con \texttt{ood\_beta=8.0}
vuelve más decisivo el cierre de la compuerta cuando el error VQ normalizado
supera el umbral. Un \texttt{familiarity\_width} de 0.5 estrecha la región
considerada familiar. Al bloquear la familiaridad durante fase 3, el notebook
evita que el entrenamiento posterior ensanche esa señal para permitir
intervenciones excesivas. Finalmente, \texttt{gate\_dropout=0.1},
\texttt{lr\_gate=1e-4} y la entropía sobre \texttt{gate\_raw} reducen la
tendencia del MLP del gate a colapsar hacia uno.

\section{Protocolo experimental}

El dataset se construye con \texttt{torchvision.datasets.Flowers102}. Los splits
\texttt{train} y \texttt{val} se concatenan y se vuelven a dividir con
\texttt{StratifiedShuffleSplit}, usando \texttt{val\_ratio=0.2} y
\texttt{split\_seed=42}. El split \texttt{test} oficial se conserva para la
evaluación final. Las imágenes se redimensionan a $224\times224$ y se normalizan
con medias y desviaciones ImageNet. En entrenamiento se añaden
\texttt{RandomHorizontalFlip} y \texttt{ColorJitter}.

Las corrupciones de entrenamiento se aplican con probabilidad 0.7 y se escogen
entre ruido gaussiano, máscara de píxeles, cutout y blur. Para evaluación se
define una suite estricta, determinista por generador semillado, con las
severidades de la Tabla~\ref{tab:corruptions}.

\begin{table}[h]
\centering
\begin{tabular}{ll}
\toprule
Corrupción & Severidades \\
\midrule
\texttt{gaussian\_noise} & 0.5, 1.0, 1.5 \\
\texttt{pixel\_mask} & 0.25, 0.5, 0.75 \\
\texttt{cutout} & 0.2, 0.35, 0.5 \\
\texttt{blur} & 0.35, 0.6, 0.85 \\
\bottomrule
\end{tabular}
\caption{Suite estricta de evaluación definida en el notebook.}
\label{tab:corruptions}
\end{table}

La evaluación mide precisión final, precisión del clasificador auxiliar,
intervenciones beneficiosas, intervenciones dañinas, tasa de pareidolia,
entropía del gate y estadísticas internas de \texttt{dq\_norm},
\texttt{uncertainty}, \texttt{familiarity}, \texttt{conflict},
\texttt{ood\_risk}, \texttt{gate\_prior}, \texttt{gate\_raw} y \texttt{gate}.
El criterio de aceptación estricta exige simultáneamente:
\[
\texttt{clean\_acc}\geq0.748252,\quad
\texttt{corrupt\_acc\_avg}\geq0.399889,
\]
\[
\texttt{gate\_order\_margin}\geq0.03,\quad
\texttt{harmful\_changes}\leq0.006099,
\]
\[
\texttt{pareidolia\_rate}\leq0.005922,\quad
\texttt{gate\_raw\_mean}\leq0.95
\]
en limpio, blur y cutout. El margen de orden se define como:
\[
\texttt{gate\_order\_margin}
=\min(\overline{g}_{blur},\overline{g}_{cutout})
-\overline{g}_{gaussian\ heavy}.
\]
La condición expresa la intención funcional de la memoria: abrir más para blur y
cutout, cerrar más para ruido gaussiano severo.

\section{Resultados}

El notebook no recomputa por defecto toda la suite; lee los artefactos finales
ubicados bajo
\texttt{experiments/VQ/out/artifacts/dememte\_attractor\_memory\_20260511\_115053/}.
Los archivos usados son \texttt{attractor\_memory\_results.csv},
\texttt{attractor\_signal\_curves.csv} y
\texttt{e5\_combined\_dropout\_ood\_tau\_150\_best.pt}.

\subsection{Comparación global}

La Tabla~\ref{tab:main_results} resume las métricas principales del CSV cargado
por el notebook. E5 no es el modelo con mayor precisión limpia absoluta, pero es
el único que satisface la aceptación estricta y alcanza el mayor promedio bajo
corrupciones.

\begin{table}[h]
\centering
\small
\begin{tabular}{lrrrrrrc}
\toprule
Modelo & Clean & Corrupt & Gauss & Mask & Cutout & Blur & Strict \\
\midrule
ResNet18 baseline & 0.6993 & 0.2869 & 0.0707 & 0.0434 & 0.5173 & 0.5162 & No \\
DeMemteAttractor & 0.7718 & 0.4207 & 0.2091 & 0.1219 & 0.6726 & 0.6791 & No \\
E1 freeze VQ P3 & 0.4931 & 0.2178 & 0.0816 & 0.0733 & 0.3481 & 0.3684 & No \\
E2 VQ clean-only P3 & 0.7671 & 0.3848 & 0.1348 & 0.0860 & 0.6415 & 0.6770 & No \\
E3 dropout + low LR & 0.7795 & 0.4036 & 0.1880 & 0.1151 & 0.6516 & 0.6596 & No \\
E4 OOD tau 0.75 & 0.7295 & 0.4083 & 0.2133 & 0.1412 & 0.6314 & 0.6475 & No \\
E4 OOD tau 1.00 & 0.7759 & 0.4010 & 0.1686 & 0.0849 & 0.6683 & 0.6820 & No \\
E4 OOD tau 1.50 & 0.7691 & 0.4228 & 0.2266 & 0.1513 & 0.6482 & 0.6649 & No \\
\textbf{E5} & \textbf{0.7653} & \textbf{0.4313} & \textbf{0.2471} & \textbf{0.1628} & \textbf{0.6570} & \textbf{0.6584} & \textbf{Sí} \\
\bottomrule
\end{tabular}
\caption{Resultados principales cargados desde \texttt{attractor\_memory\_results.csv}.}
\label{tab:main_results}
\end{table}

La comparación revela el intercambio central. El DeMemteAttractor base logra
0.7718 en limpio y 0.4207 en corrupto promedio, pero falla el criterio de orden
del gate y no supera la aceptación estricta. E3 y E4 corrigen parcialmente el
ordenamiento, pero todavía fallan alguna condición de aceptación. E5 conserva
precisión limpia por encima del piso, sube el promedio corrupto a 0.4313 y
cumple el conjunto completo de restricciones.

\subsection{Comportamiento de la compuerta}

La Tabla~\ref{tab:gate_results} muestra por qué E5 es cualitativamente distinta.
El modelo base tiene gate medio alrededor de 0.9 incluso en ruido gaussiano
severo; por tanto, la memoria interviene casi siempre. E5 reduce el gate limpio
a 0.2005 y el gate en gaussiano severo a 0.0887.

\begin{table}[h]
\centering
\small
\begin{tabular}{lrrrrrr}
\toprule
Modelo & Gate clean & Gate blur & Gate cutout & Gate gauss heavy & Margen & Pareidolia \\
\midrule
DeMemteAttractor & 0.8968 & 0.8935 & 0.9068 & 0.9225 & -0.0290 & 0.005380 \\
E3 dropout + low LR & 0.3755 & 0.3832 & 0.3799 & 0.3047 & 0.0752 & 0.000867 \\
E4 OOD tau 1.50 & 0.3710 & 0.3971 & 0.3639 & 0.1220 & 0.2419 & 0.001138 \\
\textbf{E5} & \textbf{0.2005} & \textbf{0.2104} & \textbf{0.1976} & \textbf{0.0887} & \textbf{0.1089} & \textbf{0.000257} \\
\bottomrule
\end{tabular}
\caption{Indicadores de gate y pareidolia en variantes clave.}
\label{tab:gate_results}
\end{table}

La reducción de pareidolia es particularmente relevante. E5 reporta
\texttt{pareidolia\_rate=0.000257}, muy por debajo del máximo permitido
0.005922. También mantiene \texttt{harmful\_changes=0.005895}, por debajo del
umbral 0.006099. Estas dos cifras muestran que el aumento de robustez no se
consigue mediante una memoria agresiva, sino mediante una intervención
selectiva.

\subsection{Curvas de severidad}

El CSV de curvas internas confirma el mecanismo. En limpio, el gate medio de E5
es 0.2005. En blur, aumenta ligeramente con la severidad hasta 0.2191; en cutout
permanece cerca de 0.20 y baja a 0.1867 en severidad 0.5. En cambio, frente a
ruido gaussiano pasa de 0.1536 a 0.1082 y finalmente a 0.0887. Frente a
\texttt{pixel\_mask} severo, cae hasta 0.0405.

\begin{longtable}{lrrrrr}
\caption{Curvas principales de E5 por corrupción y severidad.}
\label{tab:e5_curves}\\
\toprule
Condición & Severidad & Acc & Gate & Gate raw & Familiarity \\
\midrule
\endfirsthead
\toprule
Condición & Severidad & Acc & Gate & Gate raw & Familiarity \\
\midrule
\endhead
clean & 0.00 & 0.7653 & 0.2005 & 0.4604 & 0.3660 \\
gaussian\_noise & 0.50 & 0.5053 & 0.1536 & 0.4667 & 0.2718 \\
gaussian\_noise & 1.00 & 0.1857 & 0.1082 & 0.4612 & 0.1805 \\
gaussian\_noise & 1.50 & 0.0504 & 0.0887 & 0.4577 & 0.1417 \\
pixel\_mask & 0.25 & 0.3131 & 0.2035 & 0.4754 & 0.3698 \\
pixel\_mask & 0.50 & 0.1379 & 0.1410 & 0.4671 & 0.2459 \\
pixel\_mask & 0.75 & 0.0374 & 0.0405 & 0.4493 & 0.0428 \\
cutout & 0.20 & 0.7497 & 0.2029 & 0.4593 & 0.3713 \\
cutout & 0.35 & 0.6876 & 0.2031 & 0.4641 & 0.3713 \\
cutout & 0.50 & 0.5336 & 0.1867 & 0.4682 & 0.3383 \\
blur & 0.35 & 0.7367 & 0.2030 & 0.4561 & 0.3710 \\
blur & 0.60 & 0.6783 & 0.2091 & 0.4585 & 0.3831 \\
blur & 0.85 & 0.5601 & 0.2191 & 0.4700 & 0.4017 \\
\bottomrule
\end{longtable}

El patrón es consistente con la intención del diseño: la familiaridad se mantiene
razonable en blur y cutout, por lo que la memoria puede seguir aportando; en
ruido gaussiano y máscara extrema, la familiaridad cae y el gate se cierra.

\section{Discusión}

La evolución desde DeMemteAttractor base hasta E5 puede entenderse como el paso
de una memoria casi siempre activa a una memoria regulada por evidencia. El
modelo base ya demuestra que el atractor puede mejorar robustez: supera al
baseline ResNet18 tanto en limpio como bajo corrupciones. Sin embargo, sus
métricas internas revelan un problema: \texttt{gate\_raw\_mean\_clean=0.9997} y
valores equivalentes en blur y cutout. La compuerta cruda está saturada. En ese
régimen, el prior puede modular, pero la red aprendida ya no expresa duda real.

E5 corrige esa saturación. Su \texttt{gate\_raw\_mean\_clean=0.4604},
\texttt{gate\_raw\_mean\_blur=0.4615} y
\texttt{gate\_raw\_mean\_cutout=0.4639}. Estos valores quedan muy por debajo del
máximo de aceptación 0.95 y muestran que el MLP del gate permanece en una zona
dinámica. La combinación de dropout, learning rate menor y regularización de
entropía sobre \texttt{gate\_raw} permite esta desaturación.

El segundo elemento decisivo es la calibración OOD. Con \texttt{ood\_tau=1.5} y
\texttt{ood\_beta=8.0}, la señal de riesgo OOD se vuelve más selectiva. En vez
de abrir la memoria ante cualquier dificultad, E5 diferencia dificultad
recuperable de desviación no familiar. El resultado cuantitativo es el margen de
orden 0.1089: blur y cutout reciben más gate que ruido gaussiano severo.

El tercer elemento es el bloqueo de familiaridad. Al usar
\texttt{familiarity\_width=0.5} y \texttt{phase3\_lock\_familiarity=True}, el
entrenamiento final no puede resolver la presión de clasificación ensanchando
artificialmente la región familiar. Esto es importante porque una familiaridad
demasiado permisiva haría que el prior del gate permitiera intervención en
entradas poco confiables.

En conjunto, E5 no se limita a obtener una mejor cifra promedio. Su resultado es
una solución más coherente con el objetivo del proyecto: usar memoria cuando
ayuda y limitarla cuando puede causar errores inducidos. Por eso una variante
como E3, con mayor precisión limpia, no es seleccionada como final: todavía no
cumple todos los criterios estrictos. E5 gana porque equilibra precisión,
robustez y seguridad de intervención.

\section{Reproducibilidad}

El notebook fija semilla 42 para \texttt{random}, \texttt{numpy}, \texttt{torch}
y CUDA. Exige GPU explícitamente; si CUDA no está disponible, lanza un
\texttt{RuntimeError}. Además, incluye dos interruptores:
\texttt{RUN\_FULL\_EVAL=False} y \texttt{RUN\_TRAINING=False}. Con la
configuración por defecto, la libreta actúa como documento reproducible de
inspección: carga resultados y checkpoint sin ejecutar corridas largas. Si se
activan esos flags, puede recomputar la suite completa o reentrenar la variante
desde la fase 1.

La carga del checkpoint se realiza mediante \texttt{load\_e5\_checkpoint}, que
reconstruye el modelo E5 y carga el \texttt{state\_dict}. El
\texttt{sanity\_check\_forward} verifica un batch pequeño y comprueba shapes de
logits, gate, prior, gate crudo y finitud de SigReg. Este bloque garantiza que
el checkpoint guardado corresponde a la definición arquitectónica del notebook.

\section{Conclusión}

\texttt{e5\_final\_clean.ipynb} consolida una variante final de
DeMemteAttractor cuyo mérito principal es la regulación selectiva de una memoria
latente. La arquitectura combina un backbone ResNet18 congelado, cuantización
vectorial 2D, un atractor residual y una compuerta de ambigüedad basada en
incertidumbre, familiaridad, conflicto y riesgo OOD. El entrenamiento en tres
fases estabiliza primero la memoria, aprende después la interacción entre
atractor y gate, y finalmente ajusta el sistema con clasificación limpia y
corrupta más penalización anti-pareidolia.

La variante ganadora E5 surge al calibrar explícitamente el cierre del gate:
\texttt{ood\_tau=1.5}, \texttt{ood\_beta=8.0}, familiaridad más estrecha,
familiaridad bloqueada en fase 3, dropout en el gate, menor tasa de aprendizaje
del gate y regularización de entropía sobre \texttt{gate\_raw}. Esta combinación
reduce saturación, mejora el orden funcional de la compuerta y mantiene baja la
tasa de cambios dañinos.

En los resultados guardados que el notebook carga, E5 alcanza 0.7653 de
precisión limpia, 0.4313 de precisión corrupta promedio, 0.1089 de margen de
orden del gate, 0.005895 de cambios dañinos y 0.000257 de pareidolia. Es la
única variante listada con \texttt{strict\_acceptance\_success=True}. Por esa
razón, dentro del marco experimental codificado en el notebook, E5 constituye la
variante final ganadora.

\end{document}
